\chapter{The Central Limit Theorem}

There is a moment in everyone's probabilistic education when a strange and beautiful
fact arrives. You take a random quantity -- it does not matter much which one. It could
be the roll of a die, the outcome of a coin flip, the waiting time at a bus stop, the
glitch in a measurement. You add up many independent copies of it. And no matter how
lopsided, skewed, lumpy, or jagged the original quantity was, the \emph{shape} of the
sum settles down to the same smooth, symmetric bell curve. Every time. The bell appears
whether you started with coins or dice or exponential waiting times or the strangest
home-made distribution you can invent.

This is the Central Limit Theorem (CLT), and it is the closest thing probability has to
a law of nature. It is why the normal distribution is everywhere. It is why your
calculator's statistics functions, your A/B test, your political poll, your physics
experiment, and your machine-learning confidence band all secretly run on the same
engine. In this chapter we will state it carefully, watch it happen in simulation,
understand \emph{why} it is true (with an honest sketch using moment generating
functions), learn the practical tricks that make it usable (standardization, the
continuity correction), and find out how fast it kicks in (Berry--Esseen). By the end
you should not just believe the CLT -- you should feel that it could hardly have been
otherwise.

\section{The phenomenon, first}
\label{ch21:phenomenon}

Let us begin not with a theorem but with a picture in your head. Roll a single fair
six-sided die. The outcome is uniform on $\set{1,2,3,4,5,6}$ -- a flat, boxy
distribution, six bars of equal height. Nothing bell-shaped about it.

Now roll \emph{two} dice and look at the sum. The sum runs from $2$ to $12$, and its
distribution is no longer flat: it is a triangle, peaked at $7$, because there are more
ways to make a $7$ ($1{+}6, 2{+}5, \dots$) than to make a $2$ ($1{+}1$ only). Already
the act of summing has pulled probability toward the middle.

Roll three dice, four dice, ten dice, and the triangle rounds into a hump; the hump
gets smoother and more symmetric; and if you rescale it so it does not just slide off to
the right and spread out forever, it converges to a single fixed curve -- the standard
normal. The die was just the starting point. The summing did the work.

\begin{intuitionbox}
The CLT is a statement about \emph{averaging out idiosyncrasy}. Any one random draw
carries the full personality of its distribution -- its skew, its discreteness, its
fat tail. But when you add many independent draws, each one's personality is one voice
in a large crowd. The crowd's aggregate behavior is governed not by any individual
quirk but by just two summary numbers: the mean (where the crowd centers) and the
variance (how spread out it is). Everything else washes out. The normal distribution
is precisely the distribution that remembers \emph{only} a mean and a variance.
\end{intuitionbox}

Figure~\ref{ch21:fig-dice} shows the dice story numerically. Watch the boxy uniform
become a triangle become a bell.

\begin{figure}[H]
\centering
\begin{tikzpicture}
\begin{axis}[
  width=13cm, height=5.2cm,
  ybar, bar width=3pt,
  xlabel={value of the sum}, ylabel={probability},
  xmin=0, xmax=37, ymin=0,
  title={Sum of $n$ fair dice: $n=1$ (flat), $n=2$ (triangle), $n=5$ (bell)},
  legend style={font=\footnotesize, at={(0.98,0.97)}, anchor=north east},
]
% n=1: uniform on 1..6, prob 1/6 ~ 0.1667
\addplot+[draw=sourceblue, fill=sourceblue!40] coordinates {
 (1,0.1667)(2,0.1667)(3,0.1667)(4,0.1667)(5,0.1667)(6,0.1667)};
\addlegendentry{$n=1$}
% n=2: triangular, sum 2..12, counts/36
\addplot+[draw=quotientorange, fill=quotientorange!40] coordinates {
 (2,0.0278)(3,0.0556)(4,0.0833)(5,0.1111)(6,0.1389)(7,0.1667)
 (8,0.1389)(9,0.1111)(10,0.0833)(11,0.0556)(12,0.0278)};
\addlegendentry{$n=2$}
% n=5: sum 5..30, exact counts/7776, symmetric about mean 17.5
\addplot+[draw=densityred, fill=densityred!40] coordinates {
 (5,0.00013)(6,0.00064)(7,0.00193)(8,0.00450)(9,0.00900)(10,0.01620)
 (11,0.02636)(12,0.03922)(13,0.05401)(14,0.06944)(15,0.08372)
 (16,0.09452)(17,0.10031)(18,0.10031)(19,0.09452)(20,0.08372)
 (21,0.06944)(22,0.05401)(23,0.03922)(24,0.02636)(25,0.01620)
 (26,0.00900)(27,0.00450)(28,0.00193)(29,0.00064)(30,0.00013)};
\addlegendentry{$n=5$}
\end{axis}
\end{tikzpicture}
\caption{The distribution of the sum of $n$ fair dice. With $n=1$ it is uniform (flat).
With $n=2$ it is a perfect triangle. With $n=5$ it is already strikingly bell-shaped,
even though the underlying die has nothing bell-shaped about it.}
\label{ch21:fig-dice}
\end{figure}

\begin{historybox}
The bell curve was discovered repeatedly by people chasing exactly this phenomenon.
Abraham de Moivre, around 1733, was trying to approximate binomial probabilities --
sums of coin flips -- for large numbers of trials, because the factorials were
impossible to compute by hand. He found that the binomial probabilities, suitably
scaled, hugged a smooth curve involving $e^{-x^2/2}$. This was the very first central
limit theorem, the \emph{de Moivre--Laplace theorem}, though nobody called it that yet.
Laplace generalized it dramatically around 1810, and over the nineteenth century
Chebyshev, Markov, and Lyapunov turned it into the rigorous statement we use today.
The name ``central'' is due to George P\'olya in 1920; he meant central in the sense of
\emph{centrally important to probability theory}, not central as in middle. He was
right.
\end{historybox}

\section{Standardization: stopping the curve from running away}
\label{ch21:standardization}

Before we can state the theorem we have to deal with a bookkeeping problem. If
$X_1, X_2, \dots$ are independent and identically distributed (i.i.d.) with mean $\mu$
and variance $\sigma^2$, then the raw sum
\[
  S_n = X_1 + X_2 + \cdots + X_n
\]
has mean $n\mu$ and variance $n\sigma^2$. Both of those grow without bound. The center
of $S_n$ marches off to infinity at rate $n$, and its spread widens at rate $\sqrt{n}$.
So $S_n$ itself has no limiting distribution at all -- it just drifts away and fattens.
We have to chase it with a coordinate change.

There are two natural fixes, and they are the same fix in disguise.

\textbf{Fix the center: the sample mean.} Divide by $n$:
\[
  \bar X_n = \frac{S_n}{n}, \qquad \E[\bar X_n] = \mu, \qquad \Var(\bar X_n) = \frac{\sigma^2}{n}.
\]
Now the center is pinned at $\mu$, but the variance shrinks to zero. By the Law of Large
Numbers $\bar X_n \to \mu$. The mean stops running away, but it collapses to a point --
it loses all its shape. We have over-corrected.

\textbf{Fix both: standardize.} Subtract the mean, divide by the standard deviation.
Define
\[
  Z_n = \frac{S_n - n\mu}{\sigma\sqrt{n}}
      = \frac{\bar X_n - \mu}{\sigma/\sqrt{n}}.
\]
By construction $\E[Z_n] = 0$ and $\Var(Z_n) = 1$ for every $n$. We have removed the
drift and frozen the spread. Now $Z_n$ is neither running away nor collapsing; it is
sitting still, mean $0$ and variance $1$, ready to converge to something with a fixed
shape. That something is the standard normal.

\begin{focusbox}
The magic factor is $\sqrt{n}$ in the denominator. The sum's fluctuations around its
mean are of size $\sigma\sqrt{n}$ -- not $n$, not $\sqrt[3]{n}$, but exactly $\sqrt{n}$.
This is the single most important number in all of statistics. It says: to halve your
uncertainty about an average, you must \emph{quadruple} your sample size. Polls,
experiments, and Monte Carlo simulations all pay this $\sqrt{n}$ tax.
\end{focusbox}

\begin{whybox}
Why $\sqrt{n}$ and not something else? Variance adds for independent quantities:
$\Var(S_n) = \Var(X_1) + \cdots + \Var(X_n) = n\sigma^2$. Standard deviation is the
square root of variance, so $\SD(S_n) = \sigma\sqrt{n}$. The independence is doing the
work. If the $X_i$ were perfectly correlated (all equal), the sum would just be $nX_1$,
with standard deviation $n\sigma$ -- linear, not square-root. The $\sqrt{n}$ savings is
the dividend that independence pays you. Half of statistics is the harvest of that
dividend; the other half is worrying about whether independence really holds.
\end{whybox}

\section{The statement}
\label{ch21:statement}

\begin{theorem}[Central Limit Theorem, i.i.d. version]
\label{ch21:thm-clt}
Let $X_1, X_2, \dots$ be i.i.d.\ random variables with finite mean $\mu = \E[X_1]$ and
finite, nonzero variance $\sigma^2 = \Var(X_1)$. Let $S_n = \sum_{i=1}^n X_i$. Then the
standardized sum converges in distribution to a standard normal:
\[
  Z_n = \frac{S_n - n\mu}{\sigma\sqrt{n}} \;\dto\; \Normal(0,1).
\]
Equivalently, for every real number $z$,
\[
  \P\!\left( \frac{S_n - n\mu}{\sigma\sqrt{n}} \le z \right) \longrightarrow \Phi(z)
  \quad \text{as } n \to \infty,
\]
where $\Phi(z) = \int_{-\infty}^{z} \frac{1}{\sqrt{2\pi}} e^{-t^2/2}\,\dd t$ is the
standard normal cumulative distribution function.
\end{theorem}

Read that carefully, because the precision is where the meaning lives.

\begin{itemize}
\item \textbf{The hypotheses are minimal.} We need i.i.d., a finite mean, and a finite
nonzero variance. We need \emph{nothing} about the shape of the $X_i$. They can be
discrete or continuous, symmetric or skewed, bounded or unbounded. As long as the
variance is finite, the bell appears.
\item \textbf{The conclusion is universal.} The limit is \emph{always} the same standard
normal $\Normal(0,1)$, no matter what the $X_i$ were. The original distribution's
identity survives only through the two numbers $\mu$ and $\sigma^2$ used to standardize.
\item \textbf{Convergence is in distribution.} This means the CDFs converge,
$\P(Z_n \le z) \to \Phi(z)$, pointwise at every $z$ (here, every $z$, because $\Phi$ is
continuous). It does \emph{not} say $Z_n$ converges to any particular random number. It
is a statement about probabilities and shapes, not about realized values.
\end{itemize}

\begin{trapbox}
The finite-variance hypothesis is not a technicality you can wave away. If the $X_i$ have
infinite variance -- like a Cauchy distribution, whose tails are so heavy that
$\E[X^2] = \infty$ and even $\E[X]$ fails to exist -- the CLT is \emph{false}. The
average of $n$ independent Cauchy variables is not approximately normal; it is exactly
Cauchy again, no matter how large $n$ is. Averaging does not help at all. Heavy tails
break the theorem because a single freak observation can dominate the whole sum, so the
``crowd of equal voices'' picture collapses. There is a beautiful generalization (the
\emph{stable laws}, with $\alpha$-stable limits) for these cases, but the friendly normal
limit requires finite variance.
\end{trapbox}

In the language of the sample mean, dividing top and bottom by $n$ gives the form you
will actually use in statistics:
\[
  \frac{\bar X_n - \mu}{\sigma/\sqrt{n}} \;\dto\; \Normal(0,1),
  \qquad \text{informally} \qquad
  \bar X_n \;\approx\; \Normal\!\left(\mu, \frac{\sigma^2}{n}\right).
\]
The sample mean of a large sample is approximately normal, centered at the true mean,
with standard deviation $\sigma/\sqrt{n}$ -- the famous \emph{standard error}.

\section{Watching it happen}
\label{ch21:simulation}

The fastest way to believe the CLT is to simulate it. Take a deliberately ugly starting
distribution and watch its standardized sums smooth into a bell. We will use a very
skewed one: the exponential distribution, which has a hard wall at $0$ and a long right
tail. It looks nothing like a bell.

Figure~\ref{ch21:fig-converge} overlays the density of the standardized sum
$Z_n = (S_n - n\mu)/(\sigma\sqrt n)$ for an exponential starting variable, at $n=1$,
$n=3$, and $n=30$, against the target standard-normal curve. At $n=1$ it is the raw
exponential shape (skewed, capped on the left). By $n=3$ the skew is visibly draining
away. By $n=30$ it is nearly indistinguishable from the bell.

\begin{figure}[H]
\centering
\begin{tikzpicture}
\begin{axis}[
  width=13cm, height=6cm,
  xlabel={standardized value $z$}, ylabel={density},
  xmin=-3.5, xmax=4.2, ymin=0, ymax=0.55,
  title={Standardized sum of $n$ exponentials converging to $\Normal(0,1)$},
  legend style={font=\footnotesize, at={(0.98,0.97)}, anchor=north east},
  domain=-3.5:4.2, samples=180,
]
% target standard normal
\addplot[densityred, very thick] {0.3989422804*exp(-x^2/2)};
\addlegendentry{$\Normal(0,1)$ target}
% n=1: standardized Exp(1): density of Z = X-1 for x>=-1 is e^{-(z+1)}
\addplot[sourceblue, thick, domain=-1:4.2] {exp(-(x+1))};
\addlegendentry{$n=1$ (raw exponential)}
% n=3: standardized Gamma(3,1): z=(s-3)/sqrt3, s=sqrt3 z +3, jacobian sqrt3
\addplot[quotientorange, thick, domain=-1.732:4.2, samples=160]
  {1.732050808*((1.732050808*x+3)^2)*exp(-(1.732050808*x+3))/2};
\addlegendentry{$n=3$}
% n=30: standardized Gamma(30,1) is essentially the bell; plot a near-Gaussian
% (tiny residual right-shift in mode; variance ~1.01, peak 1/sqrt(2 pi*1.01)~0.3970)
\addplot[targetgreen, thick, dashed, domain=-3.5:4.2]
  {0.3970*exp(-(x-0.09)^2/2.02)};
\addlegendentry{$n=30$ (essentially the bell)}
\end{axis}
\end{tikzpicture}
\caption{The exponential distribution is strongly right-skewed, yet its standardized
sums march steadily toward the standard normal. At $n=1$ the curve is the raw skewed
exponential; by $n=3$ most of the skew is gone; by $n=30$ it is visually the bell. The
left wall of the exponential is what takes longest to erode.}
\label{ch21:fig-converge}
\end{figure}

\begin{happeningbox}
Look at what is physically happening as $n$ grows in
Figure~\ref{ch21:fig-converge}. The skewed exponential has a hard floor: a single draw
can never be more than $1$ standard deviation below its mean (since the minimum value
is $0$ and the mean is $1$). When you standardize a single draw, nothing can land left
of $z=-1$. But the \emph{sum} of many draws can dip below its mean by many standard
deviations, because the floor on each individual term gets averaged into irrelevance.
So the left tail, which was amputated at $n=1$, slowly grows back as $n$ increases until
it matches the normal's symmetric left tail. The CLT is the story of a distribution's
constraints dissolving under aggregation. The slowest part to heal is always the
skewness; that is exactly what the Berry--Esseen bound in
Section~\ref{ch21:berry-esseen} will quantify.
\end{happeningbox}

\begin{labbox}
Try this yourself with a few lines of code (pseudocode; any language works):
\begin{enumerate}
\item Pick a wild distribution: maybe \texttt{X = (a uniform)\^{}4}, or a 90/10 mixture
of a spike at $0$ and a uniform on $[0,10]$. Compute its true $\mu$ and $\sigma$
(estimate them from a million draws if you must).
\item For each of $n \in \set{1, 2, 5, 30, 100}$: simulate $S_n$ a hundred thousand
times, standardize each replicate to $z = (S_n - n\mu)/(\sigma\sqrt n)$, and histogram
the $z$ values.
\item Overlay the standard normal density. Watch the histogram crawl toward the bell as
$n$ grows.
\end{enumerate}
Then deliberately break it: use a Cauchy or a Pareto with index $\le 2$ (infinite
variance) and watch the histogram \emph{refuse} to converge. Seeing the failure case is
as instructive as seeing the success.
\end{labbox}

\section{Why it is true: the moment generating function sketch}
\label{ch21:why}

The simulations make the CLT believable. Now let us see \emph{why} it must be true. The
cleanest honest argument uses the moment generating function (MGF). It is not fully
rigorous -- the airtight proof uses characteristic functions, which are MGFs with an
imaginary argument and which always exist -- but the MGF version shows you the whole
mechanism with ordinary calculus, and the upgrade to characteristic functions is purely
technical.

Recall the MGF of a random variable $Y$ is $M_Y(t) = \E[e^{tY}]$, when it exists in a
neighborhood of $t=0$. Two facts make MGFs the perfect tool here.

\begin{itemize}
\item \textbf{MGFs turn sums into products.} If $Y_1, \dots, Y_n$ are independent, then
$M_{Y_1 + \cdots + Y_n}(t) = \prod_i M_{Y_i}(t)$, because the expectation of a product
of independent things factors. Summing random variables, which is hard, becomes
multiplying their MGFs, which is easy.
\item \textbf{MGFs determine distributions, and convergence of MGFs implies
convergence in distribution.} This is the \emph{continuity theorem}: if
$M_{Z_n}(t) \to M_Z(t)$ for all $t$ near $0$, and $M_Z$ is the MGF of some distribution,
then $Z_n \dto Z$. So if we can show the standardized sum's MGF converges to the normal's
MGF, we are done.
\end{itemize}

The standard normal's MGF is the target we are aiming at:
\[
  M_{\Normal(0,1)}(t) = \E[e^{tZ}] = e^{t^2/2}.
\]
Our whole job is to show $M_{Z_n}(t) \to e^{t^2/2}$.

\textbf{Setup.} Without loss of generality center and scale the $X_i$ so that $\mu = 0$
and $\sigma = 1$ (just replace $X_i$ by $(X_i - \mu)/\sigma$; this changes nothing about
the shape of the limit). Then $S_n = X_1 + \cdots + X_n$ and
$Z_n = S_n / \sqrt{n}$. Let $M(t) = \E[e^{tX_1}]$ be the common MGF of the centered,
scaled variables. Because $\E[X_1] = 0$ and $\E[X_1^2] = 1$, the Taylor expansion of $M$
around $t=0$ begins
\[
  M(t) = \E\!\left[ 1 + tX_1 + \tfrac{t^2}{2} X_1^2 + \cdots \right]
       = 1 + t\cdot 0 + \tfrac{t^2}{2}\cdot 1 + o(t^2)
       = 1 + \tfrac{t^2}{2} + o(t^2).
\]
This is the heart of the matter: \emph{the first two derivatives of $M$ at $0$ are
exactly the mean and the second moment}, and after centering and scaling those are $0$
and $1$. The personality of $X_1$ -- everything beyond its mean and variance -- has been
banished into the $o(t^2)$ remainder, and that remainder is about to be crushed.

\textbf{Multiply, because the terms are independent.} Each $X_i / \sqrt n$ has MGF
$M(t/\sqrt n)$, so
\[
  M_{Z_n}(t) = \E\!\left[ \exp\!\Big( \tfrac{t}{\sqrt n}\textstyle\sum_i X_i \Big) \right]
             = \prod_{i=1}^n M\!\left( \tfrac{t}{\sqrt n} \right)
             = \left[ M\!\left( \tfrac{t}{\sqrt n} \right) \right]^{n}.
\]
Now substitute the Taylor expansion with the small argument $t/\sqrt n$:
\[
  M\!\left( \tfrac{t}{\sqrt n} \right)
   = 1 + \frac{1}{2}\!\left(\frac{t}{\sqrt n}\right)^{\!2} + o\!\left(\frac{1}{n}\right)
   = 1 + \frac{t^2}{2n} + o\!\left(\frac{1}{n}\right).
\]
Raise to the $n$-th power and take the limit. This is the classic
$\left(1 + \frac{c}{n}\right)^n \to e^{c}$ structure, with $c = t^2/2$:
\[
  M_{Z_n}(t) = \left[ 1 + \frac{t^2/2}{n} + o\!\left(\tfrac1n\right) \right]^{n}
             \;\longrightarrow\; e^{t^2/2}.
\]
That is exactly the standard normal MGF. By the continuity theorem,
$Z_n \dto \Normal(0,1)$. $\qed$

\begin{whybox}
Step back and see what just happened, because it explains the \emph{universality}. The
only features of $X_1$ that survived into the limit were the coefficients of $t^0$,
$t^1$, and $t^2$ in the Taylor expansion of $M$ -- that is, $1$, the mean, and the
second moment. We forced the mean to $0$ and the variance to $1$ by standardizing. So
every legal starting distribution produces the \emph{identical} expansion
$1 + t^2/(2n) + \cdots$. The higher moments (skewness lives in the $t^3$ term, kurtosis
in the $t^4$ term) all entered with extra powers of $1/\sqrt n$ and got annihilated in
the limit. The $\big(1 + \tfrac{c}{n}\big)^n \to e^c$ machinery does not care where $c$
came from. \emph{That} is why the limit is always the same bell: the limit only ever
sees a mean and a variance.
\end{whybox}

\begin{connectionbox}
The limit $\big(1 + \frac{c}{n}\big)^n \to e^c$ is the same one that defines the number
$e$ and that drives continuous compound interest in calculus -- see the companion
\emph{Calculus From the Ground Up} for the careful $\eps$--$N$ version. It is striking
that the Central Limit Theorem, the law of measurement error, and the exponential
function all rest on the same little limit. The function $e^{t^2/2}$ is the unique MGF
fixed by ``mean $0$, variance $1$, and nothing else,'' which is another way of saying the
normal is the unique distribution remembering only two moments. The full proof replaces
$M(t)=\E[e^{tX}]$ with the characteristic function $\varphi(t)=\E[e^{itX}]$; since
$\abs{e^{itX}}=1$ always, $\varphi$ exists for \emph{every} distribution, even ones with
infinite MGF, so the same Taylor-then-$(1+c/n)^n$ argument goes through with no
existence worries. The complex analysis is bookkeeping; the idea above is the whole
story.
\end{connectionbox}

\section{The de Moivre--Laplace special case: coins}
\label{ch21:demoivre}

The oldest and most concrete CLT is for coin flips. Let $X_i$ be i.i.d.\ $\Bern(p)$, so
$S_n \sim \Bin(n, p)$ counts successes in $n$ trials. Here $\mu = p$ and
$\sigma^2 = p(1-p)$, so the CLT specializes to
\[
  \frac{S_n - np}{\sqrt{n p (1-p)}} \;\dto\; \Normal(0,1).
\]
This is the de Moivre--Laplace theorem, and it is enormously useful because binomial
probabilities are tedious to compute exactly for large $n$ but trivial to approximate
with a normal.

\begin{example}[A poll]
\label{ch21:ex-poll}
You poll $n = 1000$ voters; suppose the true support for a candidate is $p = 0.52$. What
is the probability your sample shows a \emph{minority} -- that is, $\bar X_n \le 0.50$?

Here $\mu = 0.52$, $\sigma = \sqrt{0.52 \cdot 0.48} = \sqrt{0.2496} \approx 0.4996$, and
the standard error of the mean is $\sigma/\sqrt n = 0.4996/\sqrt{1000} \approx 0.0158$.
Standardize the threshold:
\[
  z = \frac{0.50 - 0.52}{0.0158} \approx -1.266.
\]
So $\P(\bar X_n \le 0.50) \approx \Phi(-1.266) \approx 0.103$. Even with a true majority,
roughly a one-in-ten poll of $1000$ people would show the candidate behind. Want to cut
that to one in a hundred? You would need $z \approx -2.33$, which means shrinking the
standard error by a factor of $2.33/1.266 \approx 1.84$, which by the $\sqrt n$ law
requires about $1.84^2 \approx 3.4$ times the sample, i.e.\ around $3400$ voters. The
square-root tax in action.
\end{example}

\begin{figure}[H]
\centering
\begin{tikzpicture}
\begin{axis}[
  width=13cm, height=5.5cm,
  ybar, bar width=5pt,
  xlabel={number of heads}, ylabel={probability},
  xmin=12, xmax=38, ymin=0, ymax=0.13,
  title={$\Bin(50,0.5)$ bars vs.\ its normal approximation},
  legend style={font=\footnotesize, at={(0.02,0.97)}, anchor=north west},
]
% Binomial(50,0.5) pmf near the center (approx values)
\addplot+[draw=sourceblue!60, fill=histfill!70] coordinates {
 (15,0.0020)(17,0.0087)(19,0.0270)(21,0.0598)(23,0.0960)
 (25,0.1123)(27,0.0960)(29,0.0598)(31,0.0270)(33,0.0087)(35,0.0020)};
\addlegendentry{$\Bin(50,0.5)$ pmf (even values)}
% Normal(25, 12.5): density 1/sqrt(2*pi*12.5) e^{-(x-25)^2/25}
\addplot[densityred, very thick, domain=12:38, samples=120]
  {0.112838*exp(-(x-25)^2/25)};
\addlegendentry{$\Normal(25,12.5)$}
\end{axis}
\end{tikzpicture}
\caption{The binomial $\Bin(50, 0.5)$ probability mass function (bars, even values
shown for clarity) with its normal approximation $\Normal(25, 12.5)$ (smooth curve) laid
over it. The continuous curve threads through the tops of the discrete bars; to convert
between them honestly we need the continuity correction of
Section~\ref{ch21:continuity}.}
\label{ch21:fig-binom}
\end{figure}

\section{The continuity correction}
\label{ch21:continuity}

There is a subtlety when you approximate a \emph{discrete} distribution (like the
binomial, which lives on the integers) by a \emph{continuous} one (the normal). The
binomial puts a lump of probability exactly at each integer $k$; the normal smears
probability continuously and assigns probability $0$ to any single point. To translate
faithfully, you must decide how to split the gaps between integers.

The fix is the \emph{continuity correction}: represent the integer $k$ by the interval
$[k - \tfrac12, k + \tfrac12]$ when you pass to the normal. Each integer ``owns'' the
half-unit on either side of it. Concretely:
\[
  \P(S_n = k) \approx \P\!\left(k - \tfrac12 \le \Normal \le k + \tfrac12\right),
\]
\[
  \P(S_n \le k) \approx \Phi\!\left(\frac{k + \tfrac12 - n\mu}{\sigma\sqrt n}\right),
  \qquad
  \P(S_n \ge k) \approx 1 - \Phi\!\left(\frac{k - \tfrac12 - n\mu}{\sigma\sqrt n}\right).
\]
The half-unit shifts are easy to get backwards, so reason about which integers you mean
to \emph{include}. ``At most $k$'' includes $k$, so push the boundary up to $k+\tfrac12$.
``At least $k$'' includes $k$, so push the boundary down to $k-\tfrac12$.

\begin{example}[Continuity correction pays off]
\label{ch21:ex-cc}
Flip a fair coin $n = 100$ times. What is $\P(S_n \le 45)$? Here $\mu = 0.5$,
$\sigma = 0.5$, so $n\mu = 50$ and $\sigma\sqrt n = 0.5 \cdot 10 = 5$.

Without the correction:
$\Phi\!\big(\frac{45 - 50}{5}\big) = \Phi(-1.0) \approx 0.1587$.

With the correction:
$\Phi\!\big(\frac{45.5 - 50}{5}\big) = \Phi(-0.9) \approx 0.1841$.

The exact binomial answer is $0.1841$ to four places. The corrected approximation nails
it; the uncorrected one is off by nearly three percentage points. The smaller your $n$
or the closer you are to the tails, the more the half-unit matters.
\end{example}

\begin{trapbox}
Two common slips. First, forgetting the correction entirely when the problem is plainly
discrete -- this is the single most frequent CLT error on exams. Second, applying it to
a genuinely continuous quantity, where it does nothing but add noise: if your $X_i$ are
continuous (waiting times, measurements), there are no integer lumps and no correction
to make. The continuity correction is for discrete-to-continuous translation only.
\end{trapbox}

\section{How fast? The Berry--Esseen theorem}
\label{ch21:berry-esseen}

The CLT says $\P(Z_n \le z) \to \Phi(z)$, but it is silent about \emph{how fast}. For a
nearly symmetric, light-tailed starting distribution, $n = 20$ may already give an
excellent fit; for a violently skewed one, even $n = 200$ may leave a visible error in
the tails. We would like a guarantee. Berry (1941) and Esseen (1942) provided one.

\begin{theorem}[Berry--Esseen]
\label{ch21:thm-be}
Let $X_1, \dots, X_n$ be i.i.d.\ with mean $\mu$, variance $\sigma^2 > 0$, and finite
third absolute moment $\rho = \E\big[\,\abs{X_1 - \mu}^3\,\big] < \infty$. Let $F_n$ be
the CDF of the standardized sum $Z_n$. Then there is a universal constant $C$ such that
for all $n$ and all $z$,
\[
  \sup_{z \in \R} \big| F_n(z) - \Phi(z) \big| \;\le\; \frac{C\,\rho}{\sigma^3 \sqrt n}.
\]
The best known constant is $C < 0.4748$.
\end{theorem}

Two things to read off this beautiful inequality.

\textbf{The rate is $1/\sqrt n$.} The worst-case error between the true CDF and the
normal CDF shrinks like $1/\sqrt n$ -- the same $\sqrt n$ that governs everything else
in this chapter. Quadruple the sample, halve the approximation error.

\textbf{The constant scales with skewness.} The factor $\rho/\sigma^3$ is essentially
the absolute skewness of the starting distribution. A symmetric, concentrated
distribution has small $\rho/\sigma^3$ and converges fast; a lopsided, heavy-shouldered
one has large $\rho/\sigma^3$ and converges slowly. This is the precise version of the
intuition from the exponential simulation: \emph{skewness is what takes longest to wash
out}.

\begin{example}[Berry--Esseen for a coin]
\label{ch21:ex-be}
For a fair coin, $X_1 \in \set{0,1}$ with $\mu = \tfrac12$, $\sigma = \tfrac12$, and
$\rho = \E\,\abs{X_1 - \tfrac12}^3 = (\tfrac12)^3 = \tfrac18$. So
$\rho/\sigma^3 = (1/8)/(1/8) = 1$, and the bound reads
$\sup_z |F_n - \Phi| \le 0.4748/\sqrt n$. At $n = 100$ that is about $0.047$ -- a
worst-case CDF error under five percentage points, and in practice (away from the
extreme tails) far better, which is why the continuity-corrected answer in
Example~\ref{ch21:ex-cc} was essentially exact.
\end{example}

\begin{aestheticbox}
Berry--Esseen turns the CLT from a promise into a contract. The plain CLT is an
asymptotic statement -- true ``in the limit,'' which by itself never tells you whether
$n = 30$ is enough. Berry--Esseen gives a concrete, computable error bar for finite $n$,
depending only on three numbers you can estimate from data ($\sigma$, $\rho$, and $n$).
This is the difference between ``the bridge will eventually be safe'' and ``the bridge
holds 40 tons with this margin.'' Modern statistics and machine learning live on
finite-sample guarantees, and Berry--Esseen is the prototype.
\end{aestheticbox}

\section{Why the normal, and not some other curve?}
\label{ch21:universality}

It is worth pausing on the deepest question the CLT raises: of all possible shapes, why
is \emph{this} one -- $e^{-x^2/2}$ -- the universal attractor? Three complementary
answers.

\textbf{The fixed-point answer.} Suppose the limit of standardized sums exists and call
its law $L$. If you take two independent copies of $L$ and add them (rescaling to keep
variance $1$), you must get $L$ back -- because a sum of two big sums is just a bigger
sum, and the limit cannot depend on how you grouped the terms. So $L$ is a
\emph{fixed point} of the ``add two independent copies and rescale'' operation. The
normal is exactly such a fixed point: if $A, B \sim \Normal(0,1)$ independently, then
$(A+B)/\sqrt 2 \sim \Normal(0,1)$ again. The bell is the shape that reproduces itself
under averaging, and the CLT says it is the stable one that everything flows toward.

\textbf{The maximum-entropy answer.} Among all distributions with a given mean and
variance, the normal is the one of maximum entropy -- the most spread-out, least
committal, ``most random'' distribution consistent with those two constraints. Summing
independent things destroys information about the individual terms while preserving the
total mean and variance, so the result drifts toward the most-uncertain distribution
compatible with what survives. The normal is where information about everything-but-the-
mean-and-variance goes to die.

\textbf{The moment answer.} We saw it in the proof: the only moments that survive
standardization and the $n \to \infty$ limit are the mean and variance. Any limit law is
therefore pinned down by mean $0$ and variance $1$ alone, and the normal is the
distribution determined by exactly those. There is nothing left for the limit to be.

\begin{paradoxbox}
\textbf{The paradox of the disappearing details.} It feels almost paradoxical that the
limit forgets so much. Two starting distributions can be wildly different -- a coin and
an exponential and a uniform-to-the-fourth-power -- yet their standardized sums all
become \emph{the identical} curve. Where did all that distinguishing information go? It
did not vanish from any finite sum: for finite $n$, the sum of coins and the sum of
exponentials really are different distributions, distinguishable by their skewness and
higher moments. What the CLT says is that those differences shrink at rate $1/\sqrt n$
(skewness) and faster (kurtosis), while the bell-shaped part stays $O(1)$. So in the
\emph{ratio}, the personality becomes negligible. The information is not destroyed; it is
out-grown. This is also a warning: never conclude that two processes are ``the same''
because their averages look normal. They agree only to leading order, and the
disagreements live in the tails -- precisely where rare, important events hide.
\end{paradoxbox}

\section{Applications: the normal everywhere}
\label{ch21:applications}

\begin{applicationbox}
\textbf{Why measurement error is normal.} A single physical measurement is corrupted by
a great many tiny, independent perturbations -- thermal jitter, air currents, electronic
noise, rounding, vibration. Each contributes a small random nudge; the total error is
their sum. By the CLT, that sum is approximately normal, with mean $0$ if the
perturbations are unbiased. This is why Gauss introduced the normal distribution
precisely to model astronomical measurement error (it is still sometimes called the
Gaussian), and why ``assume the noise is Gaussian'' is the default in physics,
engineering, and signal processing. The assumption is not laziness; it is the CLT.
\end{applicationbox}

\begin{applicationbox}
\textbf{Confidence intervals and the standard error.} The single most-used consequence
of the CLT in all of data analysis. Given a sample of size $n$ with sample mean
$\bar X_n$ and sample standard deviation $s$, the CLT (plus a little extra work to
replace $\sigma$ by $s$) gives the approximate $95\%$ confidence interval
\[
  \bar X_n \pm 1.96 \cdot \frac{s}{\sqrt n}.
\]
The $1.96$ is the $z$-value cutting off $2.5\%$ in each normal tail; the $s/\sqrt n$ is
the standard error. Every poll's ``margin of error,'' every A/B test's significance
calculation, and every reported ``mean $\pm$ standard error'' in a scientific paper is
this formula, powered by the CLT. We develop it fully in the estimation and
hypothesis-testing chapters; here just notice that it requires \emph{nothing} about the
shape of the underlying data, only finite variance and a decent sample size.
\end{applicationbox}

\begin{applicationbox}
\textbf{A/B testing.} You show version A of a webpage to $n_A$ users and version B to
$n_B$, and measure conversion rates $\hat p_A$ and $\hat p_B$. Each rate is an average
of Bernoulli outcomes, so by de Moivre--Laplace each is approximately normal. Their
difference is then approximately normal too, with standard error
$\sqrt{\hat p_A(1-\hat p_A)/n_A + \hat p_B(1-\hat p_B)/n_B}$. Dividing the observed
difference by this standard error gives a $z$-statistic you can compare to the normal
curve to decide whether B genuinely beats A or just got lucky. The entire online
experimentation industry runs on the CLT applied to coin flips.
\end{applicationbox}

\begin{applicationbox}
\textbf{Monte Carlo simulation.} To estimate an integral or an expectation
$\theta = \E[g(X)]$ that is too hard to compute analytically, draw $n$ independent
samples and average: $\hat\theta_n = \frac1n \sum_i g(X_i)$. The Law of Large Numbers
says $\hat\theta_n \to \theta$; the CLT tells you the \emph{error}: $\hat\theta_n$ is
approximately $\Normal(\theta, \tau^2/n)$, so the Monte Carlo error shrinks like
$\tau/\sqrt n$. This $1/\sqrt n$ rate is independent of the dimension of $X$ -- which is
exactly why Monte Carlo beats grid-based numerical integration in high dimensions, and
why it is the workhorse of computational physics, Bayesian statistics, and quantitative
finance. See the Monte Carlo chapter for variance-reduction tricks that shrink $\tau$
(you cannot beat the $\sqrt n$, so you attack the constant).
\end{applicationbox}

\begin{connectionbox}
\textbf{Where the CLT sits among the limit theorems.} The Law of Large Numbers (previous
chapter) and the Central Limit Theorem are a matched pair, and the difference is the
scaling. The LLN looks at $\bar X_n - \mu$ on its natural scale and finds it goes to
$0$: the average converges to the truth. The CLT zooms in by the factor $\sqrt n$ and
asks what the residual fluctuation $\sqrt n\,(\bar X_n - \mu)$ looks like under the
microscope -- and finds a normal. So the LLN is the \emph{first-order} fact (the average
is right) and the CLT is the \emph{second-order} fact (here is the shape of how wrong it
is). The modes-of-convergence chapter makes ``converges in distribution'' precise and
relates it to the stronger convergences the LLN uses.
\end{connectionbox}

\section{A worked example, end to end}
\label{ch21:worked}

\begin{example}[Total weight on an elevator]
\label{ch21:ex-elevator}
An elevator is rated for $1600$ kg. Suppose adult passenger weights are i.i.d.\ with
mean $\mu = 75$ kg and standard deviation $\sigma = 15$ kg (the distribution of weights
is skewed, but we never need its shape). If $20$ people board, what is the probability
the total weight exceeds the limit?

The total $S_{20}$ has mean $20 \cdot 75 = 1500$ kg and standard deviation
$\sigma\sqrt{n} = 15\sqrt{20} \approx 67.08$ kg. We want $\P(S_{20} > 1600)$.
Standardize:
\[
  z = \frac{1600 - 1500}{67.08} \approx 1.49.
\]
Weights are continuous, so no continuity correction is needed. Then
\[
  \P(S_{20} > 1600) \approx 1 - \Phi(1.49) \approx 1 - 0.9319 = 0.068.
\]
About a $6.8\%$ chance of overload. Notice the full power of the CLT here: we knew only
the mean and standard deviation of an individual weight -- not the distribution's shape
at all -- and yet we produced a usable probability for the total. That is the everyday
miracle the theorem performs. (How trustworthy is the $6.8\%$? With $n=20$ and a
moderately skewed weight distribution, Berry--Esseen warns the answer could be off by a
few percentage points in the tail; for a safety-critical decision you would want a
larger margin or the real weight distribution.)
\end{example}

Figure~\ref{ch21:fig-tail} shows what the calculation looks like geometrically: the
total weight is approximately normal, and the overload probability is the area in the
right tail beyond the standardized threshold $z \approx 1.49$.

\begin{figure}[H]
\centering
\begin{tikzpicture}
\begin{axis}[
  width=13cm, height=5cm,
  xlabel={standardized total weight $z$}, ylabel={density},
  xmin=-3.6, xmax=3.6, ymin=0, ymax=0.45,
  domain=-3.6:3.6, samples=160,
  ytick=\empty,
  axis lines=left,
]
% shaded right tail beyond z=1.49
\addplot[densityred, very thick, name path=bell] {0.3989422804*exp(-x^2/2)};
\addplot[draw=none, name path=axis] {0};
\addplot[densityred!25] fill between[of=bell and axis, soft clip={domain=1.49:3.6}];
\draw[dashed, gray] (axis cs:1.49,0) -- (axis cs:1.49,0.13);
\node[anchor=south, font=\footnotesize] at (axis cs:2.25,0.04)
  {$\approx 0.068$};
\node[anchor=south, font=\footnotesize] at (axis cs:1.49,0.135)
  {$z\approx 1.49$};
\end{axis}
\end{tikzpicture}
\caption{The elevator overload probability as a normal tail area. The total weight of
20 passengers is approximately $\Normal(1500, 67.08^2)$; the $1600$ kg limit sits
$z \approx 1.49$ standard deviations above the mean, and the shaded right tail has area
$1 - \Phi(1.49) \approx 0.068$.}
\label{ch21:fig-tail}
\end{figure}

\section{Summary}
\label{ch21:summary}

\begin{focusbox}
\textbf{The chapter in one breath.} Add many independent finite-variance things,
subtract the mean $n\mu$, divide by $\sigma\sqrt n$, and you always get the same bell:
\[
  \frac{S_n - n\mu}{\sigma\sqrt n} \dto \Normal(0,1).
\]
\begin{itemize}
\item \textbf{Standardization} keeps the curve from drifting ($-n\mu$) or collapsing
($\div\,\sigma\sqrt n$). The $\sqrt n$ is the master scale of statistics.
\item \textbf{Why it works:} MGFs turn sums into products; Taylor expansion leaves only
the mean and variance; $(1 + \frac{t^2/2}{n})^n \to e^{t^2/2}$, the normal MGF. Higher
moments die in the limit -- hence universality.
\item \textbf{Continuity correction:} for discrete sums, replace integer $k$ by
$[k \pm \tfrac12]$ before going to the normal.
\item \textbf{Berry--Esseen:} the CDF error is at most $0.4748\,\rho/(\sigma^3\sqrt n)$
-- rate $1/\sqrt n$, scaled by skewness.
\item \textbf{Caveat:} finite variance is required. Infinite-variance (Cauchy,
heavy-tail) distributions have non-normal stable limits.
\end{itemize}
\end{focusbox}

\section*{Exercises}
\addcontentsline{toc}{section}{Exercises}

\begin{exercise}[Standardize a sum]
\label{ch21:ex-standardize}
Let $X_1, \dots, X_{40}$ be i.i.d.\ with mean $\mu = 3$ and variance $\sigma^2 = 4$. Let
$S = \sum_i X_i$. (a) Find $\E[S]$ and $\SD(S)$. (b) Using the CLT, estimate
$\P(S > 140)$. (c) Estimate $\P(115 \le S \le 130)$.
\end{exercise}

\begin{exercise}[Sample size and the square-root law]
\label{ch21:ex-sqrtn}
A quantity has standard deviation $\sigma = 8$. You want to estimate its mean with a
$95\%$ confidence interval of total width at most $1$ (i.e.\ $\pm 0.5$). Using the
interval $\bar X_n \pm 1.96\,\sigma/\sqrt n$, find the smallest $n$. Then find the $n$
needed for width $0.5$. By what factor did $n$ grow, and why?
\end{exercise}

\begin{exercise}[Continuity correction]
\label{ch21:ex-cc-prob}
A fair die is rolled $180$ times. Let $K$ be the number of sixes. (a) Find $\E[K]$ and
$\SD(K)$. (b) Use the normal approximation \emph{with} the continuity correction to
estimate $\P(K \ge 36)$. (c) Redo (b) \emph{without} the correction and report the
difference.
\end{exercise}

\begin{exercise}[A poll]
\label{ch21:ex-poll2}
In a population, $40\%$ favor a proposal. You sample $n = 600$ people. (a) What is the
approximate distribution of the sample proportion $\hat p$? (b) Estimate
$\P(\hat p \ge 0.43)$. (c) How large would $n$ need to be so that the standard error of
$\hat p$ is at most $0.01$?
\end{exercise}

\begin{exercise}[Berry--Esseen by hand]
\label{ch21:ex-be-hand}
Let $X_i$ be i.i.d.\ $\Bern(0.1)$. (a) Compute $\mu$, $\sigma^2$, and
$\rho = \E\,\abs{X_1 - \mu}^3$. (b) Compute the Berry--Esseen bound
$0.4748\,\rho/(\sigma^3\sqrt n)$ for $n = 100$ and $n = 10000$. (c) Compare to the bound
for $\Bern(0.5)$ from Example~\ref{ch21:ex-be}. Which converges to normal faster, and
does that match the intuition about skewness?
\end{exercise}

\begin{exercise}[Where the CLT fails]
\label{ch21:ex-cauchy}
The average of $n$ independent standard Cauchy random variables is again standard Cauchy,
for every $n$. (a) Explain why this does \emph{not} contradict the CLT. (b) The Cauchy
density is $f(x) = \frac{1}{\pi(1+x^2)}$. Show that $\E[X]$ does not exist by examining
$\int \abs{x} f(x)\,\dd x$. (c) In one sentence, what feature of the Cauchy breaks the
hypotheses of Theorem~\ref{ch21:thm-clt}?
\end{exercise}

\begin{exercise}[Sum of uniforms]
\label{ch21:ex-irwin}
Let $U_1, \dots, U_{12}$ be i.i.d.\ $\Unif(0,1)$. (a) Find the mean and variance of
$S = \sum_{i=1}^{12} U_i$. (b) Show that $S - 6$ has mean $0$ and variance $1$. (This is
the classic ``cheap Gaussian generator'': $\sum_{i=1}^{12} U_i - 6$ is approximately
$\Normal(0,1)$.) (c) Name one way this approximation differs from a true standard normal.
\end{exercise}

\begin{exercise}[Proof: standardization arithmetic]
\label{ch21:ex-proof-std}
Let $X_1, \dots, X_n$ be i.i.d.\ with mean $\mu$ and variance $\sigma^2$. Using only
linearity of expectation and the additivity of variance for independent variables, prove
carefully that $Z_n = (S_n - n\mu)/(\sigma\sqrt n)$ has $\E[Z_n] = 0$ and
$\Var(Z_n) = 1$ for every $n$. State exactly where independence is used.
\end{exercise}

\begin{exercise}[Proof: the MGF of a sum]
\label{ch21:ex-proof-mgf}
Let $Y_1, \dots, Y_n$ be independent with MGFs $M_{Y_i}(t)$ existing near $0$. (a) Prove
$M_{Y_1 + \cdots + Y_n}(t) = \prod_i M_{Y_i}(t)$, naming where independence enters.
(b) Specialize to the i.i.d.\ case to get $M_{S_n}(t) = M(t)^n$. (c) Using the
substitution $t \mapsto t/(\sigma\sqrt n)$, write down $M_{Z_n}(t)$ in terms of $M$ for
the centered case $\mu = 0$.
\end{exercise}

\begin{exercise}[Two-sample difference]
\label{ch21:ex-twosample}
An A/B test gives version A a conversion rate $\hat p_A = 0.118$ over $n_A = 5000$ users
and version B $\hat p_B = 0.131$ over $n_B = 5000$ users. (a) Estimate the standard error
of $\hat p_B - \hat p_A$. (b) Compute the $z$-statistic
$(\hat p_B - \hat p_A)/\text{SE}$. (c) Roughly, is the difference convincing at the
two-tailed $5\%$ level (i.e.\ is $\abs{z} > 1.96$)? Explain what the CLT is doing in
this calculation.
\end{exercise}
